
\documentclass[12pt]{article}
\usepackage{amsmath}
\usepackage{amsfonts}
\usepackage{amssymb}
\usepackage{geometry}
\geometry{a4paper, margin=0.8in}
\pagenumbering{gobble}

\title{Homework solutions for 02YELMA \\ Chapter: Electric Fields in Matter}
\begin{document}
\date{}
\maketitle
\vspace{-0.8in}

\begin{enumerate}

\item 

\begin{enumerate}
    \item The bound volume charge density is
    \[
    \rho_b = -\nabla \cdot \mathbf{P}.
    \]
    In cylindrical coordinates (radial symmetry),
    \[
    \nabla \cdot \mathbf{P} = \frac{1}{r}\frac{d}{dr}(r P_r).
    \]
    Since $P_r = \dfrac{a}{r}$,
    \[
    rP_r = a \quad \Rightarrow \quad \frac{d}{dr}(a) = 0.
    \]
    Therefore,
    \[
    \boxed{\rho_b = 0 \quad \text{for } 0<r\le R.}
    \]

    \item The bound surface charge density is
    \[
    \sigma_b = \mathbf{P} \cdot \hat{n}.
    \]
    At $r=R$, $\hat{n} = \hat{r}$, so
    \[
    \boxed{\sigma_b = \frac{a}{R}.}
    \]

    \item Total bound charge per unit length:
    \[
    Q_{\text{vol}} = \int \rho_b\, dV = 0.
    \]
    Surface contribution (per unit length):
    \[
    Q_{\text{surf}} = \sigma_b \cdot (2\pi R) = \frac{a}{R}(2\pi R) = 2\pi a.
    \]
    Hence,
    \[
    \boxed{\frac{Q_{\text{total}}}{\ell} = 2\pi a.}
    \]

    \item This is not contradictory: although $\rho_b=0$ everywhere in the bulk, charge can reside on the boundary. The divergence theorem allows nonzero flux (surface charge) even if the divergence vanishes locally.

    \item The divergence being zero only applies in the region $0<r<R$. At the boundary, $\mathbf{P}$ terminates abruptly, producing a surface charge. Mathematically, this corresponds to a singular (delta-function-like) contribution at $r=R$.
\end{enumerate}







\item 

\begin{enumerate}
    \item
    \[
    \rho_b = -\frac{dP_x}{dx} = -\frac{d}{dx}(kx^2) = -2kx.
    \]
    \[
    \boxed{\rho_b(x) = -2kx.}
    \]

    \item Surface charge densities:
    \[
    \sigma_b = \mathbf{P}\cdot \hat{n}.
    \]
    At $x=+L$, $\hat{n}=+\hat{x}$:
    \[
    \boxed{\sigma_b(+L) = kL^2.}
    \]
    At $x=-L$, $\hat{n}=-\hat{x}$:
    \[
    \boxed{\sigma_b(-L) = kL^2.}
    \]

    \item Volume contribution per unit area:
    \[
    Q_{\text{vol}}/A = \int_{-L}^{L} (-2kx)\,dx = 0.
    \]
    Surface contribution:
    \[
    Q_{\text{surf}}/A = kL^2 + kL^2 = 2kL^2.
    \]
    Hence,
    \[
    \boxed{Q_{\text{total}}/A = 2kL^2.}
    \]

    \item The system is \emph{not} neutral; there is net positive bound charge.

    \item Since $\rho_b(x) = -2kx$, it is negative for $x>0$ and positive for $x<0$, increasing linearly in magnitude away from the center. Thus, charge density is larger in magnitude near the boundaries. Physically, this is because polarization grows with $x^2$, so its spatial variation is strongest away from the center.
\end{enumerate}

\item 

\begin{enumerate}
    \item
    \[
    \rho_b = -\nabla \cdot \mathbf{P}.
    \]
    In spherical coordinates:
    \[
    \nabla \cdot (kr\hat{r}) = \frac{1}{r^2}\frac{d}{dr}(r^2 kr) = \frac{1}{r^2}\frac{d}{dr}(kr^3) = 3k.
    \]
    Hence,
    \[
    \boxed{\rho_b = -3k.}
    \]
    Surface charge:
    \[
    \boxed{\sigma_b = kR.}
    \]

    \item By symmetry, $\mathbf{E} = E(r)\hat{r}$.

    Using Gauss's law inside:
    \[
    Q_{\text{enc}} = \rho_b \frac{4}{3}\pi r^3 = (-3k)\frac{4}{3}\pi r^3 = -4\pi k r^3.
    \]
    \[
    E(4\pi r^2) = \frac{Q_{\text{enc}}}{\varepsilon_0} \Rightarrow
    \boxed{E_{\text{in}} = -\frac{k}{\varepsilon_0}r.}
    \]

    Outside:
    \[
    Q_{\text{tot}} = (-3k)\frac{4}{3}\pi R^3 + (kR)(4\pi R^2) = 0,
    \]
    so
    \[
    \boxed{E_{\text{out}} = 0.}
    \]

    \item Potential at the center:
    \[
    V(0) = \frac{1}{4\pi\varepsilon_0}\left(
    \int \frac{\rho_b}{r'}\,dV' + \int \frac{\sigma_b}{R}\,dA'
    \right).
    \]

    \item Volume contribution:
    \[
    V_{\text{vol}} = \frac{1}{4\pi\varepsilon_0}(-3k)\int_0^R \frac{4\pi r'^2}{r'}dr'
    = \frac{-3k}{\varepsilon_0}\int_0^R r'\,dr'
    = \frac{-3k}{\varepsilon_0}\frac{R^2}{2}.
    \]

    Surface contribution:
    \[
    V_{\text{surf}} = \frac{1}{4\pi\varepsilon_0}\frac{kR}{R}(4\pi R^2)
    = \frac{k}{\varepsilon_0}R^2.
    \]

    \item Total:
    \[
    V(0) = -\frac{3kR^2}{2\varepsilon_0} + \frac{kR^2}{\varepsilon_0}
    = \boxed{-\frac{kR^2}{2\varepsilon_0}.}
    \]

    \item The contributions do not cancel because the spatial weighting differs: volume charge is distributed throughout the sphere with $1/r'$ weighting, while surface charge is all at $r=R$. Even though total charge is zero, the potential need not vanish.
\end{enumerate}

\end{enumerate}

\end{document}
